% 流程节点间箭头均连接左右中心；下方为各阶段尚未排除的错误。
\begin{tikzpicture}[font=\figurefont\small,
 block/.style={draw=BookTeal,fill=BookTeal!8,minimum width=28mm,
   minimum height=12mm,align=center,inner sep=1.5mm},
 flow/.style={-{Stealth[length=2mm]},draw=BookInk,line width=0.9pt},
 note/.style={align=center,text width=28mm,font=\figurefont\footnotesize,
   text=BookMuted}]
 \node[block] (global) at (0,0) {全局描述\\筛选候选};
 \node[block] (local) at (3.7,0) {局部对应\\匹配细节};
 \node[block] (geo) at (7.4,0) {几何支持\\核验重排};
 \draw[flow] (global.east)--(local.west);
 \draw[flow] (local.east)--(geo.west);
 \node[note] at (0,-1.6) {可能只同类\\并非同一实例};
 \node[note] at (3.7,-1.6) {重复窗格\\可能相互错配};
 \node[note] at (7.4,-1.6) {几何匹配仍需\\满足查询定义};
 \foreach \x in {0,3.7,7.4}
   \draw[draw=BookRule,line width=0.7pt] (\x,-0.65)--(\x,-1);
\end{tikzpicture}
